\begin{longtable}{|p{1.8cm}|p{4cm}|p{7cm}|}
\caption{Kebutuhan Non-Fungsional \textit{Front End}} \\
\hline
\textbf{Kode} & \textbf{Kebutuhan} & \textbf{Deskripsi} \\ \hline
\endfirsthead

\caption[]{Kebutuhan Non-Fungsional \textit{Front End} (lanjutan)} \\
\hline
\textbf{Kode} & \textbf{Kebutuhan} & \textbf{Deskripsi} \\ \hline
\endhead

\hline
\multicolumn{3}{r}{\textit{Bersambung ke halaman berikutnya}} \\
\endfoot

\hline
\endlastfoot


NF-01 & \textit{Usability} &
\textit{Front-end} harus memiliki antarmuka yang mudah digunakan, mudah dipahami, serta mendukung akses cepat terhadap fitur utama aplikasi pada perangkat yang dimiliki pengguna. \\ \hline

NF-02 & \textit{Security} &
\textit{Front-end} harus menerapkan autentikasi, pembatasan akses berdasarkan peran pengguna, serta perlindungan terhadap data sensitif pengguna. \\ \hline

NF-03 & \textit{Efficiency} &
\textit{Front-end} harus mampu menampilkan data dan memuat halaman secara cepat, mengoptimalkan penggunaan jaringan, serta meminimalkan konsumsi daya baterai perangkat mengingat durasi perlombaan yang panjang. \\ \hline

NF-04 & \textit{Correctness} &
Informasi yang ditampilkan seperti posisi pelari, ETA, status \textit{checkpoint}, dan data operasional harus sesuai dengan data terbaru dari sistem \textit{backend}. \\ \hline

NF-05 & \textit{Reliability} &
\textit{Front-end} harus tetap dapat digunakan secara stabil selama perlombaan berlangsung dan mampu menangani gangguan jaringan atau GPS sementara. \\ \hline

NF-06 & \textit{Maintainability} &
Struktur kode \textit{front-end} harus modular, mudah dipelihara, serta memiliki dokumentasi yang memudahkan proses pengembangan lanjutan. \\ \hline

NF-07 & \textit{Testability} &
\textit{Front-end} harus mendukung proses pengujian antarmuka dan integrasi API yang memadai. \\ \hline

NF-08 & \textit{Flexibility} &
\textit{Front-end} harus mampu mendukung penyesuaian tampilan berdasarkan kebutuhan pengguna yang berbeda. \\ \hline

NF-09 & \textit{Reusability} &
Komponen antarmuka harus dapat digunakan kembali pada berbagai halaman dan fitur untuk menjaga konsistensi desain aplikasi. \\ \hline

NF-10 & \textit{Portability} &
\textit{Front-end} harus dapat berjalan dengan baik pada berbagai perangkat. \\ \hline

NF-11 & \textit{Interoperability} &
\textit{Front-end} harus mampu berintegrasi dengan layanan eksternal seperti API \textit{tracking}, autentikasi, dan layanan peta digital. \\ \hline

NF-12 & \textit{Accessibility / Environmental Fit} &
\textit{Front-end} harus tetap dapat dibaca dan digunakan dengan baik pada kondisi pencahayaan rendah (malam hari) serta saat pengguna dalam kondisi bergerak atau lelah secara fisik selama perlombaan berlangsung. \\ \hline

\end{longtable}